\section*{Capítulo \ref{ch:g11:deriv} --- Derivación: un primer curso}

\begin{solution}{exo:g11:deriv:1}
Para $f(x) = x^2$ en $a = 2$:
\[
\frac{(2+h)^2 - 4}{h} = \frac{4h + h^2}{h} = 4 + h
\xrightarrow[h\to0]{} 4,
\]
luego $f'(2) = 4$. Para $g(x) = x^2 + 3x$ en $a = 1$: $g(1) = 4$ y
\[
\frac{(1+h)^2 + 3(1+h) - 4}{h} = \frac{5h + h^2}{h} = 5 + h
\xrightarrow[h\to0]{} 5,
\]
luego $g'(1) = 5$.
\end{solution}

\begin{solution}{exo:g11:deriv:2}
$f'(x) = 12x^2 - 4x + 1$.

$g'(x) = \frac{5}{2\sqrt x} - \frac{3}{x^2}$ (en $\intoo{0}{+\infty}$).

$h'(x) = 2(x^2 - 3) + (2x+1)(2x) = 2x^2 - 6 + 4x^2 + 2x
= 6x^2 + 2x - 6$ (regla del producto).
\end{solution}

\begin{solution}{exo:g11:deriv:3}
$f(x) = x^2 - 3x$, $a = 2$: $f(2) = -2$ y $f'(x) = 2x - 3$, luego
$f'(2) = 1$; \omterm{def:g11:deriv:tangent}{tangente} $y = -2 + (x - 2) = x - 4$.

$f(x) = \frac1x$, $a = 1$: $f(1) = 1$, $f'(1) = -1$; \omterm{def:g11:deriv:tangent}{tangente}
$y = 1 - (x - 1) = 2 - x$.

$f(x) = \sqrt x$, $a = 4$: $f(4) = 2$,
$f'(4) = \frac{1}{2\sqrt4} = \frac14$; \omterm{def:g11:deriv:tangent}{tangente}
$y = 2 + \frac14(x - 4) = \frac{x}{4} + 1$.
\end{solution}

\begin{solution}{exo:g11:deriv:4}
$f'(x) = \frac{(x-1) - (x+2)}{(x-1)^2} = \frac{-3}{(x-1)^2}$.

$g'(x) = \frac{2x(x^2+1) - x^2 \times 2x}{(x^2+1)^2}
= \frac{2x}{(x^2+1)^2}$.

$h'(x) = -\frac{2x+1}{(x^2+x+1)^2}$ (regla del cociente con $u = 1$).
\end{solution}

\begin{solution}{exo:g11:deriv:5}
$f'(x) = 3x^2 - 12x + 9 = 3(x^2 - 4x + 3) = 3(x-1)(x-3)$: positiva fuera
de $\intcc{1}{3}$ y negativa dentro. Así pues, $f$ es \omterm{def:g11:func:monotone}{creciente} en
$\intoc{-\infty}{1}$, \omterm{def:g10:functions:variations}{decreciente} en $\intcc{1}{3}$ y \omterm{def:g11:func:monotone}{creciente} en
$\intco{3}{+\infty}$, con un \omterm{def:g10:functions:extrema}{máximo} relativo $f(1) = 1 - 6 + 9 - 2 = 2$ y
un \omterm{def:g10:functions:extrema}{mínimo} relativo $f(3) = 27 - 54 + 27 - 2 = -2$.
\end{solution}

\begin{solution}{exo:g11:deriv:6}
En $\intoo{-1}{+\infty}$:
\[
f'(x) = \frac{2x(x+1) - (x^2+3)}{(x+1)^2}
= \frac{x^2 + 2x - 3}{(x+1)^2}
= \frac{(x+3)(x-1)}{(x+1)^2}.
\]
Para $x > -1$, tanto $x + 3$ como $(x+1)^2$ son positivos, así que
$f'(x)$ tiene el signo de $x - 1$: $f$ es \omterm{def:g10:functions:variations}{decreciente} en $\intoc{-1}{1}$ y
\omterm{def:g11:func:monotone}{creciente} en $\intco{1}{+\infty}$, con \omterm{def:g10:functions:extrema}{mínimo} $f(1) = \frac{4}{2} = 2$.
\end{solution}

\begin{solution}{exo:g11:deriv:7}
$f'(x) = 2x - 4$.

\emph{1.} \omterm{def:g11:deriv:tangent}{Tangente} horizontal: $f'(x) = 0$ en $x = 2$; el punto es
$(2, f(2)) = (2, 1)$ (el vértice de la \omterm{def:g11:quad:parabola}{parábola}).

\emph{2.} Paralela a $y = 2x + 1$ significa \omterm{def:g10:reffunc:affine}{pendiente} $2$: $2x - 4 = 2$,
luego $x = 3$, lo que da el punto $(3, f(3)) = (3, 2)$.
\end{solution}

\begin{solution}{exo:g11:deriv:8}
Con anchura $x$ (dos lados vallados) y largo $60 - 2x$:
$A(x) = x(60 - 2x) = 60x - 2x^2$ en $\intoo{0}{30}$. Entonces
$A'(x) = 60 - 4x$, positiva antes de $x = 15$ y negativa después: el área
es máxima para $x = 15$, con un cercado de $15 \times 30$ m y área
$450$ m$^2$. La fórmula del vértice
$\alpha = -\frac{60}{2 \times (-2)} = 15$ del \cref{ch:g11:quad} da la
misma respuesta sin cálculo diferencial.
\end{solution}

\begin{solution}{exo:g11:deriv:9}
Sea $f(x) = \frac{x+1}{2} - \sqrt x$ en $\intoo{0}{+\infty}$. Entonces
\[
f'(x) = \frac12 - \frac{1}{2\sqrt x} = \frac{\sqrt x - 1}{2\sqrt x},
\]
que tiene el signo de $\sqrt x - 1$: negativa en $\intoo{0}{1}$ y positiva
en $\intoo{1}{+\infty}$. Así que $f$ decrece y después crece, con \omterm{def:g10:functions:extrema}{mínimo}
$f(1) = 1 - 1 = 0$. Por tanto, $f(x) \geq 0$ para todo $x > 0$, es decir,
$\sqrt x \leq \frac{x+1}{2}$, con igualdad solo en $x = 1$.
\end{solution}

\begin{solution}{exo:g11:deriv:10}
De $\pi r^2 h = 500$ se obtiene $h = \frac{500}{\pi r^2}$, luego
\[
S(r) = 2\pi r^2 + 2\pi r \cdot \frac{500}{\pi r^2}
= 2\pi r^2 + \frac{1000}{r},
\qquad
S'(r) = 4\pi r - \frac{1000}{r^2}.
\]
$S'(r) = 0$ cuando $4\pi r^3 = 1000$, es decir,
$r^3 = \frac{250}{\pi}$, luego $r = \sqrt[3]{250/\pi}$ ($\approx 4.3$
cm). Como $S'(r) = \frac{4\pi r^3 - 1000}{r^2}$ es negativa antes de ese
valor y positiva después, se trata de un \omterm{def:g10:functions:extrema}{mínimo}.
\end{solution}

\begin{solution}{exo:g11:deriv:11}
La \omterm{def:g11:deriv:tangent}{tangente} a $y = x^2$ en el punto de \omterm{def:g10:coordgeom:system}{abscisa} $a$ es
$y = a^2 + 2a(x - a) = 2ax - a^2$. Pasa por $P(1, -3)$ cuando
$-3 = 2a - a^2$, es decir, $a^2 - 2a - 3 = 0$, cuyas raíces son $a = 3$ y
$a = -1$ ($\Delta = 16$). Hay, por tanto, \emph{dos} \omterm{def:g11:deriv:tangent}{tangentes} que pasan
por $P$:
\[
y = 6x - 9 \quad (a = 3),
\qquad
y = -2x - 1 \quad (a = -1).
\]
\end{solution}

\begin{solution}{exo:g11:deriv:12}
Por el \cref{ex:g11:deriv:study}, $f$ crece hasta el \omterm{def:g10:functions:extrema}{máximo} relativo
$f(-1) = 3$, decrece hasta el \omterm{def:g10:functions:extrema}{mínimo} relativo $f(1) = -1$ y vuelve a
crecer; para $x$ positivo grande (resp.\ negativo), $x^3$ domina y $f$
toma valores arbitrariamente grandes (resp.\ pequeños). Leyendo las rectas
horizontales $y = k$ sobre la \omterm{def:g10:functions:table}{tabla de variación}:
\begin{itemize}
  \item $k > 3$ o $k < -1$: la recta corta la curva una vez: $1$ solución;
  \item $k = 3$ o $k = -1$: la recta toca un punto de giro y atraviesa
        otra rama: $2$ soluciones;
  \item $-1 < k < 3$: la recta atraviesa las tres ramas: $3$ soluciones.
\end{itemize}
\end{solution}

\begin{solution}{pb:g11:deriv:1}
\textbf{1.} $f'(x) = 6x - 5$; $g'(x) = 3x^2 - 12$. Por la definición:
$\frac{(1 + h)^2 - 1}{h} = 2 + h \to 2$: la \omterm{def:g11:deriv:derivative}{derivada} de $x^2$ en $1$ es
$2$.

\textbf{2.} $g(2) = 8 - 24 = -16$ y $g'(2) = 12 - 12 = 0$: la \omterm{def:g11:deriv:tangent}{tangente} es
la recta horizontal $y = -16$. Una \omterm{def:g11:deriv:tangent}{tangente} horizontal señala un punto
crítico, aquí un \omterm{def:g10:functions:extrema}{mínimo} relativo, como confirma la \omterm{def:g10:functions:table}{tabla de variación} de
la pregunta 5.

\textbf{3.} \omterm{def:g10:reffunc:affine}{Pendiente} $2a$ en $(a, a^2)$:
$y = a^2 + 2a(x - a) = 2ax - a^2$.

\textbf{4.} En $x = 0$: $y = -a^2$: la \omterm{def:g11:deriv:tangent}{tangente} corta al eje a
profundidad $a^2$ por debajo del origen, exactamente tan \emph{abajo} de
cero como $P$ está por encima. Receta: para trazar la \omterm{def:g11:deriv:tangent}{tangente} en un punto
de altura $h$, únelo con el punto del eje situado a profundidad $h$ por
debajo del origen. (No hace falta ninguna \omterm{def:g11:deriv:derivative}{derivada} en el tablero de
dibujo.)

\textbf{5.} $g'(x) = 3(x^2 - 4)$: positiva en $\intoo{-\infty}{-2}$,
negativa en $\intoo{-2}{2}$ y positiva más allá. $g$ crece hasta un \omterm{def:g10:functions:extrema}{máximo}
relativo $g(-2) = 16$, cae hasta un \omterm{def:g10:functions:extrema}{mínimo} relativo $g(2) = -16$ y vuelve
a crecer.

\textbf{6.} $F\left(0, \frac14\right)$ y $T(0, -a^2)$:
$FT = \frac14 + a^2$.

\textbf{7.} $FP = a^2 + \frac14 = FT$: isósceles en $F$, así que los
ángulos de la base son iguales: $\widehat{FTP} = \widehat{FPT}$.

\textbf{8.} El rayo entrante en $P$ es vertical y, por tanto, paralelo a
la recta $(TF)$ (el eje $y$). La \omterm{def:g11:deriv:tangent}{tangente} $(TP)$ corta a las dos
paralelas, así que el ángulo entre el rayo y la \omterm{def:g11:deriv:tangent}{tangente} es igual a
$\widehat{FTP}$ (ángulos alternos), que a su vez es igual a
$\widehat{FPT}$, el ángulo entre $[PF]$ y la \omterm{def:g11:deriv:tangent}{tangente}. Ángulos iguales a
los dos lados de la \omterm{def:g11:deriv:tangent}{tangente} en $P$.

\textbf{9.} Por la ley de la reflexión, el rayo que llega verticalmente a
$P$ sale por la recta que forma el mismo ángulo del otro lado; por la
pregunta 8, sale por $[PF]$: pasa por el \omterm{pb:g11:quad:1}{foco}, sea cual sea $a$. Al
revés, la luz nacida en $F$ sale de cada punto del espejo paralela al eje.
Las antenas recogen y los faros iluminan: demostrado.

\textbf{10.} $a = 1$: \omterm{def:g11:deriv:tangent}{tangente} $y = 2x - 1$, $T(0, -1)$;
$FP = \sqrt{1 + \left(1 - \frac14\right)^2} = \sqrt{\frac{25}{16}} =
\frac54$; $FT = \frac14 + 1 = \frac54$: isósceles, como se prometía.

\textbf{11.} La \omterm{def:g11:deriv:tangent}{tangente} en $x_n$ es $y = f(x_n) + f'(x_n)(x - x_n)$;
corta a $y = 0$ en $x = x_n - \frac{f(x_n)}{f'(x_n)}$.

\textbf{12.} $x_1 = 1 - \frac{-1}{2} = \frac32$;
$x_2 = \frac32 - \frac{1/4}{3} = \frac{17}{12}$;
$x_3 = \frac{577}{408}$: las aproximaciones de $\sqrt2$ de Herón, o sea,
la máquina de \omterm{pb:g10:functions:1}{punto fijo} del \cref{pb:g10:functions:1}. Identidad:
$x - \frac{x^2 - 2}{2x} = \frac{2x^2 - x^2 + 2}{2x}
= \frac{x^2 + 2}{2x} = \frac12\left(x + \frac2x\right)$: la receta de
Herón \emph{es} el método de Newton aplicado a $x^2 - 2$, dos milenios
antes.

\textbf{13.} $x_1 = 5 - \frac{125 - 100}{75} = \frac{14}{3} \approx
4.6667$; $x_2 \approx 4.6667 - \frac{1.63}{65.33} \approx 4.6417$.
Calculadora: $\sqrt[3]{100} = 4.6416\ldots$: cuatro decimales correctos en
dos pasos.

\textbf{14.} $f'(2) = 0$: la \omterm{def:g11:deriv:tangent}{tangente} en el punto de partida es horizontal
(pregunta 2) y no corta nunca al eje; la fórmula divide entre cero.
Advertencia: arranca el método de Newton lejos de los puntos críticos y,
en general, suficientemente cerca de la \omterm{def:g11:quad:discriminant}{raíz} que buscas.

\textbf{15.} La bisección gana una cifra binaria por paso; Newton
aproximadamente \emph{duplica} el número de cifras correctas en cada paso
($1 \to 3 \to 6$ aquí). Cuando cada paso es caro (millones de parámetros,
navegación en tiempo real), duplicar gana a avanzar a pasitos, y por eso
la recta \omterm{def:g11:deriv:tangent}{tangente} corre dentro de tanto silicio.

\textbf{16.} $d^2 = 900 - w^2$, luego $S(w) = w(900 - w^2) = 900w - w^3$,
para $0 < w < 30$.

\textbf{17.} $S'(w) = 900 - 3w^2 = 0$ en $w = \sqrt{300} = 10\sqrt3
\approx 17.3$ cm ($S' > 0$ antes y $< 0$ después: un \omterm{def:g10:functions:extrema}{máximo}). Entonces
$d = \sqrt{600} = 10\sqrt6 \approx 24.5$ cm.

\textbf{18.} $\frac{d^2}{w^2} = \frac{600}{300} = 2$, luego
$d = w\sqrt2$: el canto es la anchura por $\sqrt2$, la razón de los
carpinteros (y la construcción del diámetro dividido en tres partes la
entrega con un compás, sin álgebra sobre el terreno).

\textbf{19.} Viga cuadrada: $w = d = \sqrt{450} \approx 21.2$ cm, con
rigidez $w^3 = 450^{3/2} \approx 9\,546$. Viga óptima:
$S = 10\sqrt3 \times 600 = 6\,000\sqrt3 \approx 10\,392$. Ganancia: un
$8.9\,\%$ más rígida a partir del mismo tronco; la \omterm{def:g11:deriv:derivative}{derivada} paga el
aserradero.

\textbf{20.} Geometría: los ángulos iguales de la \omterm{def:g11:deriv:tangent}{tangente} convirtieron un
paraboloide en un recogedor de luz. Cálculo numérico: deslizarse por las
\omterm{def:g11:deriv:tangent}{tangentes} resuelve ecuaciones duplicando cifras. Optimización: una
\omterm{def:g11:deriv:derivative}{derivada} que se anula halló el rectángulo más resistente del tronco. Un
solo objeto y tres oficios; y la convexidad del año que viene dirá,
además, a qué lado de cada \omterm{def:g11:deriv:tangent}{tangente} queda la curva.
\end{solution}
